\documentclass[a4paper,11pt]{report}
\usepackage[utf8]{inputenc}
\usepackage[T1]{fontenc}
\usepackage{fixltx2e}
\usepackage{graphicx}
\usepackage{longtable}
\usepackage{float}
\usepackage{wrapfig}
\usepackage{rotating}
\usepackage[normalem]{ulem}
\usepackage{amsmath}
\usepackage{textcomp}
\usepackage{marvosym}
\usepackage{wasysym}
\usepackage{amssymb}
\usepackage{hyperref}
\tolerance=1000
\usepackage{minted}
\usepackage[spanish]{babel}
\usepackage{amsmath}
\usepackage{amsthm}
\usepackage{dsfont}

\newtheorem{theorem}{Teorema}
\newtheorem{fact}{Proposición}
\newtheorem{lemma}{Lema}
\newtheorem{corollary}{Corolario}
\newtheorem{definition}{Definición}
\setlength{\parindent}{0pt}
\setlength{\parskip}{1em}

\usepackage{color}
\newenvironment{wording}{\setlength{\parskip}{0pt}\rule{\textwidth}{0.5em}}{~\\\rule{\textwidth}{0.5em}}

\everymath{\displaystyle}


\begin{titlepage}
	\centering
	\includegraphics[width=0.5\textwidth]{./imgs/ugr.png}\par
	\vspace{1cm}
        \rule{\textwidth}{0.3em}\hfill
        {\huge\bfseries Clasificación no balanceada: \par
        Análisis teórico y algoritmos \par}
        \rule{\textwidth}{0.3em}\hfill
	\vspace{1cm}
        {\scshape Proyecto final de Carrera\par}
        {\scshape Ingeniería informática y Matemáticas\par}
        \vfill
        {\bf Autor \par}
	{Ignacio Cordón Castillo \par}
	{\bf Tutor \par}
	{Salvador García López \par}
	\vfill
      	\includegraphics[width=0.4\textwidth]{./imgs/by-nc-sa.png}\par
	{\large \today\par}
\end{titlepage}

\begin{document}
\maketitle
\tableofcontents

\break

\chapter{Introducción}
Basado mayormente en el contenido de \cite{shwartz_understanding_ml}
\chapter*{Introducción}

  \section*{Bibliografía empleada}
  El material principal en el que se ha basado la sección de matemáticas ha sido \citeauthor{shalev}, 
  que proporciona una detallada formalización de los fundamentos matemáticos del Aprendizaje Automático. Puntualizamos además
  una serie de referencias extras usadas en la elaboración de esta parte:
  
  \begin{itemize} 
   \item La primera parte del capítulo de introducción a teoría de la probabilidad ha sido confeccionado por el autor de 
   este trabajo a partir de \citet{caratheodory}, que ofrece una demostración guiada a través de ejercicios de dicho teorema. 
   También se ha usado como apoyo para dicho capítulo el libro \citeauthor{loeve}. 
   \item Para la parte final del primer capítulo de introducción a probabilidad se ha usado \citeauthor{shalev}, pero
   mayoritariamente el contenido de Wikipedia \cite{wiki:markov}, \cite{wiki:hoeff_lemma} y \cite{wiki:hoeffding}.
   \item Para el teorema \ref{th:fundamental} sobre el teorema fundamental PAC se ha usado principalmente el material de clase
   proporcionado por \citeauthor{slfetaya} para sus clases en \textit{Weizmann Institue of Science}.
  \end{itemize}
   
   

\bibliography{references}
\bibliographystyle{IEEEtran}

\end{document}